\section{Introduction}\label{sec:introduction}

Asymptotic convexity describes the convexity that remains after finitely
many directions are discarded. The question considered here is whether
symmetric midpoint convexity can persist in a reflexive Banach space
which admits no asymptotically uniformly convex equivalent norm.

All spaces and scalars are real. For an infinite-dimensional Banach
space $Y$ with norm $N$, write $\operatorname{cof}(Y)$ for the closed
finite-codimensional subspaces of $Y$. For $t>0$, define
\begin{align}
 \overline\delta_N(t)
 &=\inf_{N(x)=1}\sup_{F\in\operatorname{cof}(Y)}
   \inf_{y\in F,\,N(y)=1}\bigl(N(x+ty)-1\bigr),
   \label{eq:auc-modulus}\\
 \widehat\delta_N(t)
 &=\inf_{N(x)=1}\sup_{F\in\operatorname{cof}(Y)}
   \inf_{y\in F,\,N(y)=1}
   \left(\frac{N(x+ty)+N(x-ty)}2-1\right).
   \label{eq:average-modulus}
\end{align}
The norm is \emph{asymptotically uniformly convex} (AUC) if
$\overline\delta_N(t)>0$ for all $t>0$, and \emph{asymptotically
midpoint uniformly convex} (AMUC) if $\widehat\delta_N(t)>0$ for all
$t>0$. Thus AUC requires a gain in every sufficiently far-out direction;
AMUC permits the two opposite directions to be tested together. A norm
$N$ is equivalent to a given norm $\|\cdot\|$ if
$a\|x\|\le N(x)\le b\|x\|$ for some $0<a\le b<\infty$ and
all $x$. A space is \emph{AUC renormable} if it admits an equivalent
AUC norm.

Dilworth, Kutzarova, Randrianarivony, Revalski, and Zhivkov introduced
AMUC and constructed an AMUC norm on $\ell_2$ which is not AUC
\cite[Definition~2.2 and Theorem~2.4]{DKRRZ2016}. Their question about
equivalence up to renorming \cite[Section~5]{DKRRZ2016} has already
been answered negatively by Baudier \cite[Theorem~1 and Corollary~1]{Baudier2026},
using the Kadets--Werner modification of the Bourgain--Rosenthal space
\cite[Theorem~2.5]{KW2004}. We give a reflexive example with an explicit
midpoint estimate.

The historical definition of AMUC uses the maximum of the two signed
norms in \eqref{eq:average-modulus}. This gives the same positivity
property. The average never exceeds the maximum. Conversely, intersect
any chosen finite-codimensional subspace with the kernel of a functional
norming $x$. Both signed norms are then at least one, so a gap for their
maximum gives half that gap for their average. Our quantitative modulus
always uses the average in \eqref{eq:average-modulus}.

The example also addresses a metric question. Define the
\emph{countably branching diamonds} by taking $D_0$ to be a single edge
and replacing every edge at each step by countably many internally
disjoint two-edge paths. Each current edge has length one; let $d_k$
be the shortest-path metric on $D_k$. An embedding of $D_k$ into a
Banach space has distortion at most $D$ if it is a map $f$ for which
some scale $s>0$ satisfies
\begin{equation}\label{eq:embedding-definition}
 s\,d_k(u,v)\le\|f(u)-f(v)\|\le Ds\,d_k(u,v)
 \qquad(u,v\in D_k).
\end{equation}
The family embeds uniformly when one finite $D$ works for all $k$.
For finite binary diamonds, Lee and Naor obtained quantitative $L_p$
distortion bounds by summing convexity losses \cite[Lemma~2.2]{LeeNaor2004},
and Johnson and Schechtman developed a self-improvement proof
\cite[Section~4]{JohnsonSchechtman2009}.
Baudier, Causey, Dilworth, Kutzarova, Randrianarivony, Schlumprecht,
and Zhang proved that an equivalent AMUC norm prevents uniform
embeddings of the countably branching diamonds
\cite[Theorem~4.1]{BaudierEtAl2017}. For reflexive spaces
with unconditional asymptotic structure, they also proved that failure
of AUC renormability implies uniform diamond embeddings
\cite[Theorem~B and Corollary~5.3]{BaudierEtAl2017}. Swift extended the embedding construction to countably branching bundle
graphs under the same reflexive, unconditional, and non-AUC-renormable
hypotheses \cite[Theorem~3.2 and Corollary~3.4]{Swift2018}. Perreau
subsequently treated duals of separable spaces with weak-star unconditional asymptotic
structure, retaining unconditional asymptotic structure in the reflexive
formulation \cite[Theorems~4.8 and~5.3]{Perreau2021}. The example below
shows that the diamond implication fails for general reflexive spaces.

More recently, Basset, Lancien, and Proch\'azka related AMUC to separated
families of midpoints and obtained the diamond obstruction even when
the maps are controlled only on specified vertex pairs
\cite[Definition~6.15 and Propositions~6.16 and~6.20]{BassetLancienProchazka2026}.
Malthaner obtained quantitative obstructions for countably branching
diamonds using a Matou\v{s}ek-type coloring and path argument
\cite[Theorems~3.5 and~3.7]{Malthaner2025}. We use the full
bi-Lipschitz convention \eqref{eq:embedding-definition} and prove the
required midpoint estimate directly.

\subsection{The example and result}

For each integer $h\ge1$, let $T_h=\mathbb N^{\le h}$ be the
rooted tree of words of length at most $h$, ordered by extension of
words. Its root is the empty word. Keep separate copies of these trees
in the forest
\[
 \mathcal F=\coprod_{h\ge1}(\{h\}\times T_h).
\]
Write $a\preceq b$ if $a$ is an ancestor of $b$ in the same component.
A \emph{segment} is a nonempty interval
$[a,b]=\{v:a\preceq v\preceq b\}$ in an ancestral chain.
For a finitely supported real function $u$ on $\mathcal F$, set
$u(S)=\sum_{v\in S}u(v)$ and define
\begin{equation}\label{eq:segment-norm}
 \rho(u)=\sup\left\{\left(\sum_{i=1}^m u(S_i)^2\right)^{1/2}:
  m\ge1,\ S_1,\ldots,S_m\text{ pairwise disjoint segments}\right\}.
\end{equation}
The Euclidean triangle inequality and singleton segments show that
$\rho$ is a norm, with $\|u\|_2\le\rho(u)\le\|u\|_1$.
Let $J$ be its completion and let $X=J^*$, with
the dual norm denoted by $\|\cdot\|$. We write $e_v\in J$ for a
coordinate unit vector. Singleton tests give $|u(v)|\le\rho(u)$,
so coordinate evaluation extends to an element $e_v^*\in X$.
The construction in \eqref{eq:segment-norm} uses James's
square-sum norm on disjoint segments \cite[p.~739]{James1974}; the
branching and component heights here differ from the classical binary
tree.


For the classical binary James tree space, Girardi proved AUC of both
the coordinate predual and the full dual
\cite[Theorems~3 and~5]{Girardi2001}. The countable branching and finite,
unbounded component heights used here lead to the following different
behavior.

\begin{theorem}\label{thm:main}
The space $X=J^*$ defined by \eqref{eq:segment-norm} is infinite
dimensional, separable, and reflexive, and has no equivalent AUC norm.
Its given norm satisfies
\[
 \widehat\delta_X(t)\ge\sqrt{1+t^2/12}-1\qquad(t>0).
\]
For every integer $k\ge0$, every embedding of $D_k$ into $X$ has
distortion at least $\sqrt{1+k/12}$.
\end{theorem}

This gives a negative answer to the reflexive diamond question recorded
as Problem~39 in \cite[author manuscript p.~388]{BaudierLancien2025}. The metric obstruction is
unchanged by equivalent renorming, although the numerical constants
refer to the specified dual norm.

\subsection{How the argument fits together}

Section~\ref{sec:forest} establishes reflexivity and the coordinate
projections used throughout the proof. A sum of coordinate functionals
along one root path has norm one, while its next coordinate increment
is a weakly null unit vector as the child varies. In
Section~\ref{sec:no-auc}, a hypothetical AUC norm forces repeated
multiplicative gains along one such path. Arbitrarily high finite trees
contradict its bounded norm.

The midpoint argument in Section~\ref{sec:paired} must control segment
tests crossing from a finite ancestral set (one closed under taking ancestors) to its
complement. Simply
adding norming tests for the two parts creates uncontrolled positive
cross terms. We add one correction at their first exit vertices. The
correction makes both resulting tests inexpensive, and its evaluations
cancel when the two endpoint norms are averaged. Lemma~\ref{lem:paired}
expresses the result as a lower bound in terms of the squared head and
tail norms.
Section~\ref{sec:general-forest} extends this estimate to coordinate
completions on general rooted forests.

Finally, the infinitely many middle vertices of a diamond supply a
separated family of approximate midpoints. Two of their finite
coordinate projections can be made arbitrarily close. Their
half-difference therefore reduces the metric problem to the head--tail
estimate. Section~\ref{sec:diamonds} turns this into a loss in squared
stretch at every replacement and sums the losses. This gives a direct
quantitative instance of the midpoint and self-improvement methods for
these graphs \cite[Section~4]{BaudierEtAl2017}.
